\section{Boundedly many roots imply linear full-support MCA}
\label{app:power-family-mca}
\label{app:bounded-root-mca}

A polynomial solution variety of quadratic degree need not produce
quadratically many exceptional labels on a received line. We prove an
$O(n)$ bound for every candidate family with a bounded number of distinct
roots, allowing both the roots and their multiplicities to vary.

Let $\mathcal D$ be a set of $n$ distinct field elements and let $f,g$ be
arbitrary words on it. For a candidate $P$ of degree at most $D$, set
\[
 S(P,z)=\{x\in\mathcal D:f(x)+zg(x)=P(x)\}.
\]
Call $(P,z)$ bad at threshold $A>D$ if $|S(P,z)|\ge A$ and no
polynomials $F_0,G_0$ of degree at most $D$ agree with $f,g$, respectively,
on all of $S(P,z)$. Equivalently, $g$ has no degree-$D$ interpolant on
this full support, since otherwise $F_0=P-zG_0$ works.
A label is bad if it has at least one bad candidate in the family.

\begin{theorem}[Bounded root count]
\label{thm:bounded-root-linear-mca}
\label{thm:power-family-linear-mca}
Fix $r\ge1$ and $0<\rho_-\le\rho_+<a\le1$. Suppose
$\rho_-\le D/n\le\rho_+$ and $A\ge an$.
There are constants $C_{r,a}$, $c_{r,a}$, and
$n_0(r,a,\rho_-)$ such that, for $n\ge n_0$, every received line has
at most $C_{r,a}n$ bad labels witnessed by polynomials with at most $r$
distinct roots over an algebraic closure, if the characteristic is zero
or $p>c_{r,a}n$. Nonzero scalar coefficients, root positions, and root
multiplicities may all vary. Including the zero polynomial costs at
most $n$ additional labels.
\end{theorem}

The constants below are explicit but very large. No general
first-order MCA theorem or improved better.codes certificate follows.
We may work over an algebraic closure when discussing roots: a
degree-$D$ interpolant on more than $D$ base-field points descends to
the base field.

\subsection{Pencils and clusters with small residual degree}

An affine polynomial pencil $P_z=U+zV$, with $\deg U,\deg V\le D$,
contains at most $n$ bad labels. Indeed a bad full support contains a
coordinate where $(f,g)\ne(U,V)$; agreement there forces $g\ne V$
and determines the unique label $z=(U-f)/(g-V)$.

Suppose a selected family of distinct bad labels has candidates $HQ$,
where $H$ is fixed with at most $s$ distinct roots, and $\deg Q\le K$.
We can take $K\le D-\deg H$. If three graph points $(z,Q)$ are not
collinear, their line-incidence determinant is a nonzero polynomial of
degree at most $K$. Their common agreement support has size at most
$K+s$, after dividing by $H$ away from its roots.
Assume $K+s\le a^3n/16$. For $L\ge4/a$, counting ordered triples
through coordinates gives a lower bound $a^3nL^3/8$: if $b_x$ counts
agreeing candidates, use
$(b_x)_3\ge\max(b_x-2,0)^3$ and convexity.
Noncollinear triples contribute at most $(a^3n/16)L^3$.
Each ordered pair determines a unique graph line, which contains at
most $n$ selected bad labels by the pencil observation. Collinear
triples therefore contribute at most $n^2L^2$. Consequently
\begin{equation}\label{eq:root-cluster-small}
 L\le16n/a^3.
\end{equation}
The same bound covers $L<4/a$. This argument needs no characteristic
hypothesis.

Set $\delta=a^3/(128r)$. Fix any set $S$ of at most $3r$ algebraic
roots and consider candidates whose total root multiplicity outside
$S$ is at most $r\delta D$. Partition their exponents on $S$ into
boxes of side $w=\lfloor\delta D\rfloor+1$. Within each box factor
the product $H$ given by its lower-corner exponents. The residual degree
is at most
\[
 |S|(w-1)+r\delta D\le4r\delta D\le a^3n/32.
\]
If $3r\le a^3n/32$, equation~\eqref{eq:root-cluster-small} applies
to each box. There are at most
\[
 B_0=(1+\lceil1/\delta\rceil)^{3r}
\]
boxes, so this cluster contains at most $Cn$ selected bad labels, where
\begin{equation}\label{eq:root-cluster-C}
 C=16B_0/a^3.
\end{equation}
The bound is uniform in $S$. Factors may be considered over the
algebraic closure; a pencil through two original base-field candidates
has base-field coefficients.

\subsection{Four candidates with private high-multiplicity roots}

\begin{lemma}\label{lem:bounded-root-gcd}
Let $P_1,\ldots,P_4$ be nonzero polynomials of degree at most $D$,
each with at most $r$ distinct roots. Suppose each $P_i$ has a root
of multiplicity at least $\delta D$ which is not a root of any other
$P_j$. For an integer $m\ge1$, put
\[
 N=\binom{m+3}{3},\qquad J=m+1,\qquad M=N-J.
\]
Assume
\[
 \delta D>\binom{N-1}{2}(4r-1),\qquad
 \operatorname{char}\F=0\quad\text{or}\quad p>mD.
\]
For two independent homogeneous linear forms $F,G$ in four variables,
the gcd of $F(P_1,\ldots,P_4)$ and $G(P_1,\ldots,P_4)$, after
removing factors supported on the union of the roots of the $P_i$,
has degree at most
\begin{equation}\label{eq:bounded-root-gcd}
 \frac{24D}{m+5}+\frac{(4r-1)(M-1)}2.
\end{equation}
\end{lemma}

\begin{proof}
The $N$ homogeneous degree-$m$ monomials in the $P_i$ are linearly
independent. For otherwise take a minimal relation with $q$ terms
$f_1,\ldots,f_q$. A ratio of two distinct monomials has an order of
absolute value at least $\delta D$ at some private root, so $q=2$
is impossible. Any $q-1$ terms have a nonzero ordinary Wronskian:
their degrees are less than $p$, and row reduction to distinct leading
degrees gives a nonzero Vandermonde leading coefficient. The same
argument holds in characteristic zero.

Let $S$ be the union of the roots of the $P_i$ and $A_0$ the common
gcd of the $f_i$. All Wronskians obtained by omitting one term are
constant multiples. Omit a term of largest degree for the global
degree bound, and a term of smallest order separately at each root for
the local bounds. Comparing them yields
\[
 \max_i\deg f_i-\deg A_0\le\binom{q-1}{2}(|S|-1).
\]
The left side is at least the rational degree of any ratio $f_i/f_j$,
and hence at least $\delta D$. Since $q\le N$ and $|S|\le4r$,
this contradicts the assumed threshold.

The degree-$m$ part of the homogeneous ideal $(F,G)$ has dimension
$M$: its quotient is a polynomial ring in two variables, of degree-$m$
dimension $J$. Evaluation in the $P_i$ is injective by the preceding
paragraph. Write $d_\alpha$ for the degrees of the evaluated monomials.
A basis of the evaluated ideal has sum of degrees at most the sum of
the largest $M$ of the $d_\alpha$. Indeed, order the formal monomials
by degree and row-reduce from highest degree; each basis element has
degree at most the degree of its distinct pivot. Thus its nonzero
Wronskian $W$ satisfies
\[
 \deg W\le\sum_\alpha d_\alpha-\binom M2.
\]

At a root $s\in S$, let $a_\alpha(s)$ be the monomial orders.
Intersecting a codimension-$J$ subspace with the span of monomials
of order at least each integer $k$ gives an adapted basis whose sum of
orders is at least
$\sum_\alpha a_\alpha(s)-J\max_\alpha a_\alpha(s)$.
The dimension of each such intersection loses at most $J$;
summing these filtration dimensions proves the bound. Cancellations
can only increase orders. A constant basis change preserves the order
of $W$, so, using
\[
 \sum_{s\in S}\sum_\alpha a_\alpha(s)=\sum_\alpha d_\alpha,
 \qquad
 \sum_{s\in S}\max_\alpha a_\alpha(s)
 \le m\sum_{i=1}^4\deg P_i,
\]
we obtain
\[
 \sum_{s\in S}\operatorname{ord}_s W
 \ge\sum_\alpha d_\alpha-Jm\sum_i\deg P_i-|S|\binom M2.
\]
Let $T$ be the gcd in the statement. It divides every evaluated ideal
element. The identity
$\operatorname{Wr}(Th_1,\ldots,Th_M)=T^M\operatorname{Wr}(h_1,\ldots,h_M)$
contributes at least $M\deg T$ additional zeros outside $S$.
Combining the degree and order bounds gives
\[
 \deg T\le\frac{4JmD}{M}+\frac{(4r-1)(M-1)}2.
\]
Since $M=m(m+1)(m+5)/6$, this is \eqref{eq:bounded-root-gcd}.
\end{proof}

This is an elementary specialization of the high-power gcd phenomenon.
A broader characteristic-zero result is Guo--Sun--Wang
\cite[Theorem 11]{guo-sun-wang-gcd}; the finite-characteristic statement
above follows directly from the Wronskian proof.

\subsection{Counting the remaining quadruples}

Select one nonzero bad candidate at each bad label. Call a root heavy
if its multiplicity is at least $\delta D$. A quadruple without a
private heavy root for every member has some member $P_i$ whose heavy
roots all lie in the union of the other three root sets. Fix the other
three candidates. Their root union has size at most $3r$, and every
root of $P_i$ outside it has multiplicity less than $\delta D$.
The cluster bound \eqref{eq:root-cluster-C} gives at most $Cn$ choices
for this selected bad label. Summing over the four possible members,
there are at most $4CnL^3$ such ordered quadruples. This includes
candidates with no heavy roots.

For every other quadruple at distinct labels, common agreement
coordinates satisfy the two independent linear forms
\begin{align*}
 &(z_2-z_3)P_1+(z_3-z_1)P_2+(z_1-z_2)P_3=0,\\
 &(z_2-z_4)P_1+(z_4-z_1)P_2+(z_1-z_2)P_4=0.
\end{align*}
Lemma~\ref{lem:bounded-root-gcd}, together with the at most $4r$
roots in the union, bounds their common support. Choose
\[
 m\ge1536/a^4-5,\qquad
 \frac{(4r-1)(M-1)}2+4r\le a^4n/64,
\]
as well as the lemma's hypotheses and $3r\le a^3n/32$.
Their common support then has size at most $a^4n/32$.
For $L\ge6/a$, fourth incidence gives
\[
 \frac{a^4nL^4}{16}
 \le\frac{a^4nL^4}{32}+4Cn^2L^3.
\]
Therefore
\begin{equation}\label{eq:bounded-root-labels}
 L\le\frac{128C}{a^4}n
   =\frac{2048}{a^7}(1+\lceil128r/a^3\rceil)^{3r}n.
\end{equation}
The small-$L$ case also satisfies this bound. At fixed positive rate,
all degree and length thresholds eventually hold, and a sufficiently
large fixed lower bound on $p/n$ ensures $p>mD$.
The zero polynomial contributes at most $n$ more labels by the same
single-coordinate argument as for a fixed pencil. This proves
Theorem~\ref{thm:bounded-root-linear-mca}.

\subsection{A quadratic-degree solution surface}

Fix $e\ge1$ in characteristic zero or $p>2e$. The nonzero polynomial
solutions of degree at most $2e$ to
\begin{equation}\label{eq:quadratic-power-ode}
 e^2P^2P'''+3e(1-e)PP'P''+(1-e)(1-2e)(P')^3=0
\end{equation}
are exactly $P=cH^e$ with $\deg H\le2$. At a root of multiplicity
$m$, the coefficient of the lowest possible term, after removing a
nonzero local scalar, is $m(m-e)(m-2e)$. Thus every root has
multiplicity $e$ or $2e$. Conversely substitution verifies every
$cH^e$; equivalently this is the denominator-cleared equation
$(P^{1/e})'''=0$. The monic base descends by uniqueness.

The projective map $[H]\mapsto[H^e]$ from the plane of quadratics is
injective and base-point free, with injective tangent maps since
$e\ne0$ in the field. Its image has degree $e^2$. Explicitly, cut by
$P(u)=P(v)$ and $P(w)=P(v)$ for three distinct points $u,v,w$.
There is no nonzero solution with $H(v)=0$, so normalize $H(v)=1$.
The values $H(u),H(w)$ independently range over the $e$-th roots of
unity, giving $e^2$ distinct reduced points by quadratic interpolation.
This fixed third-order, jet-degree-three equation has an actual solution
variety of degree $\Theta(n^2)$ at fixed positive rate, but
Theorem~\ref{thm:bounded-root-linear-mca} gives only $O(n)$ bad labels
in its large-characteristic regime.

The checkers in \path{research/bounded_root_mca/} verify unequal-degree
Wronskian bookkeeping and cluster counting, including every full support
of 1,025 bad labels in a length-4,096 example with at most two roots per
candidate. The earlier \path{research/power_family_mca/} package verifies
the quadratic-power classification and degree sections. These checks
supplement the proofs. Families with a growing number of distinct roots
and the fixed-gap quadratic lower-bound target remain open here.
